% ===================================================================
\chapter{Our Results Against Kiem's, One by One}
\label{ch:kiemcompare}

This chapter brings every comparison with Kiem together in one place. For each
of his results or claims it gives this thesis's corresponding result, says
whether the two are the same, and explains why or why not, pointing to the
measurement that shows it. Hardware figures were accounted for individually in
Chapter~\ref{ch:interpret}; they are summarised here so that the whole
comparison can be read at once.

\section{Three kinds of comparison}
\label{sec:kc_kinds}

Not every pair of numbers can be compared directly, and the most common
mistake made during this project was to compare numbers that measure different
things. Each comparison below is placed in one of three classes.

\begin{description}[leftmargin=!,labelwidth=30mm,style=nextline,font=\normalfont\bfseries]
\item[Like for like] The two numbers measure the same quantity of the same
thing, and a difference is a real difference. Example: the arithmetic
specification of DRHE; post-route LUTs of two circuits.
\item[Through a bridge] The numbers come from different data or setups, but
something built for the purpose connects them: the replica of his hardware
(Chapter~\ref{ch:interpret}) or the replication of his simulation
(Chapter~\ref{ch:replication}).
\item[Not comparable] The numbers look alike but measure different things,
and no bridge closes the gap. Example: his CR on his recordings against ours on
ColoRadar. For these the chapter explains what each number does measure.
\end{description}

\section{The DRHE algorithm, line by line}
\label{sec:kc_audit}

Kiem's Section~3.5.3 specifies the final scheme and his Section~4.3.4 the
hardware form of it. Table~\ref{tab:kc_audit} takes every element in turn and
compares it with this thesis's MATLAB model (\texttt{compress\_drhe.m}) and its
HLS design (DRHE-1).

\begin{table}[p]
\centering
\caption[DRHE audit against Kiem \S{}3.5.3]{Every element of Kiem's final DRHE
scheme (his \S{}3.5.3, \S{}3.5.4 and \S{}4.3.4) against this thesis's MATLAB model and
HLS design. The letters in the last column refer to the explanations in
Section~\ref{sec:kc_why}.}
\label{tab:kc_audit}
\scriptsize
\renewcommand{\arraystretch}{1.25}
\setlength{\tabcolsep}{3pt}
\begin{tabular}{@{}>{\raggedright}p{0.16\textwidth}>{\raggedright}p{0.26\textwidth}>{\raggedright}p{0.18\textwidth}>{\raggedright}p{0.22\textwidth}>{\raggedright\arraybackslash}p{0.10\textwidth}@{}}
\toprule
Element & Kiem & MATLAB \texttt{compress\_drhe.m} & HLS DRHE-1 & Same?\\
\midrule
Input & FX16 range FFT, first $N/2$ bins, per RX, ramp by ramp & same & same (128-bit AXIS, 4 RX) & yes\\
Cartesian to polar & $a = \sqrt{re^2+im^2}$, $\theta = \mathrm{atan2}(im, re)$ & same, double & same, single precision & yes; (b)\\
Magnitude, Eq.~3.5 & $\hat a[m] = \alpha\hat a[m-1] + (1-\alpha)a[m-1]$ & same & same & yes\\
Phase, Eq.~3.6 & $\hat\theta[m] = \beta\hat\theta[m-1] + (2-\beta)\theta[m-1] - \theta[m-2]$ & same & same & yes\\
Phase wrap & once into $[-\pi,\pi]$ (his listing, \S{}4.3.4) & same & same & yes\\
$\alpha$, $\beta$ & 0.6, 0.4 (Table 4.4) & same & same & yes\\
Polar to Cartesian & rounded to integer (\texttt{AP\_RND}) & \texttt{round} & \texttt{roundf} & yes\\
Difference, Eq.~3.7 & $X - \hat X$ in 16-bit two's complement & \texttt{wrap\_int16} & \texttt{wrap\_int16} & yes\\
Initial state & zero (\S{}3.5.4) & zero per frame & zero, cleared on ramp 0 of every frame & yes; (c)\\
State update & from the input sample & from the input & from the reconstructed sample & \textbf{changed}; (a)\\
Run-length code & removed (\S{}3.5.1) & none & none & yes\\
Size category S4 & Table 3.7, 0--15 & $\lfloor\log_2|v|\rfloor+1$, capped at 15 & same regions, by comparison & yes\\
APPEND & $v$ if $v>0$; $v + 2^{S_4} - 1$ if $v<0$ & counted & same formula & yes\\
Residual $-32768$ & no code (S4 = 15 spans $\pm[16384, 32767]$); not addressed & counted as S4 = 15 & clamped to $-32767$ & \textbf{changed}; (a)\\
Dictionary & fixed, Appendix A & same 16 lengths & same codes, bit-reversed & yes; (d)\\
Arithmetic & fixed point: \texttt{ap\_ufixed<16,16>}, \texttt{ap\_fixed<16,3>}, \texttt{ap\_fixed<32,4>} & double & single-precision float & \textbf{changed}; (b)\\
Output packing & 58-bit word per channel into 256-bit AXIS & bit count only & 512-bit accumulator, 256-bit AXIS & layout; (d)\\
Prediction lag & ramp $m-1$ (one TX) & $m-1$ & $m-1$ (DRHE-1); $m-n_{\mathrm{Tx}}$ (DRHE-$n$) & extension; (e)\\
\bottomrule
\end{tabular}
\end{table}

\subsection{Verification that the algorithm is the same}

Two independent implementations of Kiem's Section~3.5.3 --- the thesis's
\texttt{compress\_drhe.m} and the replication's \texttt{kr\_drhe\_encode.m},
written from his text without reference to the first --- produce identical
residuals and identical bit counts on the same frames
(Section~\ref{sec:kr_crosscheck}). The HLS DRHE-1 matches
\texttt{compress\_drhe.m} on all 50 ColoRadar frames to within the AXI padding
of the last output word (Section~\ref{sec:phase2}), and a third implementation,
in Python, reproduces the HLS bit count of every frame exactly
(Appendix~\ref{app:verify}). The DRHE of this thesis is Kiem's DRHE.

\subsection{The five places the hardware differs, and why}
\label{sec:kc_why}

\paragraph{(a) The $-32768$ clamp and the closed loop: a defect fix.} Kiem's
Table~3.7 gives category 15 the range $\pm[16384, 32767]$, exactly $2^{15}$
values filling the 15 APPEND bits, so a residual of $-32768$ has no code; his
text does not address it. It can occur (for example $re = -32768$ with a
prediction of zero), and when it does, the APPEND formula aliases it onto
$+32767$ (Section~\ref{sec:hole}). The HLS saturates it to $-32767$. Once a
sample has been saturated, an open-loop encoder --- one updating its state from
the original sample, as Kiem's does --- disagrees with the decoder, which can
only update from the reconstruction, and because the filters are recursive the
disagreement persists for the rest of the frame. Driving the encoder's state
from the reconstruction confines the error to one sample and one LSB. Whenever
no clamp fires the reconstruction equals the input, so the closed loop is
\emph{bit-identical} to Kiem's open loop; on ColoRadar the residuals span
$[-514, +416]$ and it never fires. Its cost is hardware only: 26 cycles of
pipeline depth in floating point, about four in fixed point
(Section~\ref{sec:acc_depth}).

\paragraph{(b) Floating-point arithmetic: a Thesis~A choice, now shown to be
unnecessary.} Floating point was chosen so that the HLS would follow the
double-precision MATLAB model exactly --- Kiem reports that his MATLAB could not
be made bit-accurate to his fixed-point HLS and needed an HLS decoder for that
reason --- and because his own future work suggested it
(Section~\ref{sec:kiem_open}). On ColoRadar the float and fixed-point
predictors give residuals whose bit counts differ by \SI{0.065}{\percent}, while
the float datapath costs most of the design's logic, every extra DSP and twice
the block RAM (Chapter~\ref{ch:interpret}). The fixed-point variant is the better
engineering choice and is recommended as future work.

\paragraph{(c) Clearing the state every frame: a clarification.} Kiem's model is
``initialized with zero''; his MATLAB decodes one frame at a time, and his
hardware text does not say what happens between frames. DRHE-1 clears the state
at the first ramp of every frame, so every frame decodes on its own and a lost
frame cannot corrupt the next. The one extra write per cycle that this costs is
what holds the design at $\mathrm{II} = 2$ (Section~\ref{sec:acc_ii}).

\paragraph{(d) Bit packing: the same bits in a different order.} The codewords
and APPEND fields have exactly Kiem's lengths, so the compressed size is his
apart from padding. The codewords are stored bit-reversed because the packer is
least-significant-bit first, and the stream is therefore not bit-compatible
with a decoder written for his 58-bit packets --- it was never meant to be.
The packer costs about 4,200 LUTs more than his (ladder steps 1--2).

\paragraph{(e) DRHE-$n$: an extension.} On a single-transmitter radar the
previous ramp is the previous observation of the same channel. On the ColoRadar
cascade it is a different transmitter eleven times in twelve. DRHE-$n$ keeps one
state per transmitter and predicts from ramp $m - n_{\mathrm{Tx}}$; with
$n_{\mathrm{Tx}} = 1$ it is DRHE-1, and therefore Kiem's algorithm.

\section{Every result, compared}
\label{sec:kc_table}

Table~\ref{tab:kc_all} lists every result and claim of Kiem's that this thesis
can speak to. The sections after it discuss the ones that need more than a
line.

\begin{table}[p]
\centering
\caption[Kiem's results against this thesis's]{Kiem's results and claims
against this thesis's. ``Class'' is the kind of comparison of
Section~\ref{sec:kc_kinds}: L like for like, B through a bridge, N not
comparable.}
\label{tab:kc_all}
\scriptsize
\renewcommand{\arraystretch}{1.3}
\setlength{\tabcolsep}{3pt}
\begin{tabular}{@{}>{\raggedright}p{0.17\textwidth}>{\raggedright}p{0.14\textwidth}>{\raggedright}p{0.19\textwidth}>{\centering}p{0.05\textwidth}>{\centering}p{0.07\textwidth}>{\raggedright\arraybackslash}p{0.30\textwidth}@{}}
\toprule
Result or claim & Kiem & This thesis & Class & Same? & Why\\
\midrule
DRHE algorithm & \S{}3.5.3 & identical residuals and bits & L & yes & Two independent implementations agree to the bit; hardware differs in five documented places (Section~\ref{sec:kc_why}).\\
DRHE CR, simulation & 3.25 (real data) & 3.30 (MATLAB, 16 RX); 3.310 (HLS, 4 RX) & N & no\,/\,coincidence & Different data, sensor, scaling and mechanism (Section~\ref{sec:kc_cr}).\\
DRHE CR on his scene & --- & 2.28 at his $-75$\,dBFS; 3.25 at $-87.9$ & B & pattern & Replication reproduces 3.25/3.24 only at a noise level 17\,dB below his stated floor (Section~\ref{sec:kr_m88}).\\
RLE adds little & 3.25 vs 3.24 & 3.25 vs 3.24 on his scene; $\leq$\SI{1.6}{\percent} on ColoRadar & B, L & yes & Holds on both datasets (Section~\ref{sec:kc_rle}).\\
Fixed dictionary costs little & 3.24 $\to$ 3.17 (incl.\ HW FFT) & \SI{0.3}{\percent} on his scene at $-87.9$; a re-derived dictionary gains 12--\SI{18}{\percent} on corrected ColoRadar & B, L & his data only & The dictionary encodes his residual statistics (Section~\ref{sec:kc_dict}).\\
DRHE CR, hardware & 3.17 & 3.310 (DRHE-1), 3.751 (DRHE-$n$) & N & no & As above, and his hardware FFT is fixed point; ours takes MATLAB's FX16 codes.\\
LPC is ``the same approach'' but less promising & not measured; 2.056 quoted from \cite{meucci2022} & Meucci--Mancuso style: 1.32--1.54 on his scene; slow-time LPC: 3.697 vs DRHE-$n$ 3.751 & B & for his LPC only & Two different algorithms share the name (Section~\ref{sec:kc_lpc}).\\
DRHE is the best of his design space & yes & yes, among lossless schemes and against BAQ/RDVLE & L & yes & The lossy schemes lose detections (Table~\ref{tab:phase1}).\\
FX16 detection cost & FN 0.59\,\%, FP 1.18\,\% & FN 3.77\,\%, FP 1.11\,\% & N & no & Different sensor, scene and processing chain (Section~\ref{sec:kc_det}).\\
Noise floor & $-70.7$\,dBFS & $-105.4$\,dBFS & N & no & Different sensor and reference; not comparable (Section~\ref{sec:kc_det}).\\
LUT & \num{45892} & \HW{drheone}{LUT} & L & no & Float datapath and packer against his build; replica \HW{ablk}{LUT} (Chapter~\ref{ch:interpret}).\\
FF & \num{9243} & \HW{drheone}{FF} & L & no & Float operator registers, 58-stage pipeline.\\
BRAM (36\,Kb tiles) & 10 & 20 & L & no & 32-bit against 16-bit state words; replica matches exactly.\\
DSP & 44 & 132 & L & no & Every DSP is a float operator; replica \HW{ablk}{DSP}.\\
II & 1 & 2 & L & no & One extra state write per cycle; step 3 of the ladder reaches 1.\\
Pipeline latency & 23 cycles (HLS est.) & 58 cycles (HLS) & L & no & 26 cycles closed loop, rest float latency.\\
Clock & 100\,MHz, grade $-1$ & \HW{drheone}{fmax}\,MHz ($-2$); \HW{ablctrlsgone}{fmax}\,MHz ($-1$) & L & yes & Both meet 100\,MHz on his grade.\\
Throughput & 11.92 Gibit/s (theoretical) & 6.40 theoretical, 5.14 measured & L & factor 2 & II 2 against 1; his figure is in binary units.\\
System latency & 0.08\,ms & not measured; 3.2\,\si{\micro\second} added by the compressor & N & --- & His figure includes FFT, DMA and interrupts.\\
Fits his device & 65\,\% of LUTs & 50\,\% of LUTs & L & fits & Post-route (Section~\ref{sec:resinterp}).\\
\bottomrule
\end{tabular}
\end{table}

\section{The compression ratio: 3.25 and 3.30 are not the same number}
\label{sec:kc_cr}

The closeness of DRHE's 3.30 on ColoRadar to Kiem's 3.25 on his recordings was
taken, early on, as confirmation that the implementation was right. The
implementation \emph{is} right (Section~\ref{sec:kc_audit}), but the match is a
coincidence: the two ratios are produced by different mechanisms.

\paragraph{ColoRadar's 3.30.} Section~\ref{sec:decomposition} decomposed it.
The capture uses 9.8 of the 16 FX16 bits, which is worth a factor of \num{1.64}
before any coding; the sparsity of a range-FFT frame is worth \num{2.13}
through the entropy coder; and DRHE's predictor, applied along a
transmitter-interleaved axis, \emph{costs} \SI{4.9}{\percent}, so the no-prediction ratio, 3.479 over
50 frames, is higher than DRHE's.

\paragraph{Kiem's 3.25.} On his single-transmitter radar the predictor runs on
the right axis, and his dictionary --- built from his own residuals --- is
optimal for residuals of 6 to 8\,LSB (Section~\ref{sec:kr_fingerprint}). His
ratio is what a matched coder gives on residuals of that size; on his own
simulated scene the replication reaches it only when the FX16 noise is about
4.6\,LSB.

\paragraph{Why they cannot be compared directly.} Beyond the different scenes,
four properties of the chain differ between the two datasets, each of which moves
the ratio on its own:

\begin{itemize}
\item \textbf{FX16 scaling.} This thesis maps a full-scale sinusoid to 32767;
Kiem's FFT core to about 8191. The same signal occupies two more bits here.
\item \textbf{ADC width and gain.} Kiem's CTRX is 14-bit (12-bit in his use
case) and his data is scaled by $16 - \mathit{bitWidth} + 16$; ColoRadar is
recorded in a 16-bit container but peaks near 1100 codes, so its effective
gain is low.
\item \textbf{Real against complex samples.} His FFT input is real and half the
spectrum is redundant; ColoRadar's is complex and the negative half is
discarded by choice.
\item \textbf{One against twelve transmitters.} His ramp axis is slow time; ours
is interleaved.
\end{itemize}

\noindent
What \emph{can} be said is this. Applied to ColoRadar, Kiem's algorithm gives
exactly DRHE-1's 3.310, because it is DRHE-1. Applied to Kiem's scene, this
thesis's DRHE gives exactly his algorithm's ratio at every noise level, because
it is his algorithm. The comparison of \emph{algorithms} is therefore settled;
the comparison of \emph{datasets} cannot be made.

\section{Run-length coding}
\label{sec:kc_rle}

Kiem removed the run-length half of the original coder after finding that it
changed the ratio from 3.25 to 3.24. The same comparison, with ideal
dictionaries in both cases, holds on both datasets available here: 3.25
against 3.24 on the replicated scene at $-87.9$\,dBFS
(Table~\ref{tab:kr_m88}), and on ColoRadar (every fifth frame, 4 RX) 3.466
against 3.457 for DRHE-1 and 4.340 against 4.273 for DRHE-$n$, a gain of at
most \SI{1.6}{\percent}. His simplification is sound for this sensor too.

\section{The fixed dictionary}
\label{sec:kc_dict}

Kiem fixed his dictionary on one recording and found the cost small. On his
simulated scene at the noise level that reproduces his ratio, the fixed
dictionary costs \SI{0.3}{\percent} (3.23 against 3.24). On ColoRadar the
answer depends on the predictor. Re-deriving the dictionary from five frames and
testing it on the other 45 raises the ratio by 3--\SI{8}{\percent} for the raw
codes and the lag-1 predictors, and by 12--\SI{18}{\percent} for the corrected
predictors, whose residuals are much smaller than the ones the dictionary was
built for (Section~\ref{sec:dictresults}). A dictionary is a
property of the data it was trained on, and this one was trained on a
different radar.

\section{Linear prediction}
\label{sec:kc_lpc}

\begin{table}[htbp]
\centering
\caption[The LPC question in one table]{Kiem's statement about LPC and what was
measured. CRs with Kiem's fixed dictionary unless stated.}
\label{tab:kc_lpc}
\small
\begin{tabular}{@{}lrrr@{}}
\toprule
Comparison & DRHE & LPC & Source\\
\midrule
\multicolumn{4}{@{}l}{\emph{LPC-Huffman of \cite{meucci2022} (fast time, raw ADC) --- the method Kiem meant}}\\
Kiem's scene, $-75$\,dBFS, $p = 10$     & 1.85 / 2.28\,(R4S4) & 1.32 & Table~\ref{tab:kr_m75}\\
Kiem's scene, $-87.9$\,dBFS, $p = 10$   & 3.23 / 3.25\,(R4S4) & 1.54 & Table~\ref{tab:kr_m88}\\
\midrule
\multicolumn{4}{@{}l}{\emph{Slow-time LPC on FX16 (this thesis)}}\\
Kiem's scene, $-75$\,dBFS               & 1.85 & 2.06 & Table~\ref{tab:kr_m75}\\
Kiem's scene, $-87.9$\,dBFS             & 3.23 & 3.38 & Table~\ref{tab:kr_m88}\\
ColoRadar, lag 1 (HLS, 50 frames)       & 3.310 & 3.385 & Table~\ref{tab:tdmhls}\\
ColoRadar, lag $n_{\mathrm{Tx}}$ (HLS)  & \textbf{3.751} & 3.697 & Table~\ref{tab:tdmhls}\\
ColoRadar, lag $n_{\mathrm{Tx}}$, re-derived dictionary & 4.207 & \textbf{4.371} & Section~\ref{sec:dictresults}\\
\bottomrule
\end{tabular}
\end{table}

\noindent
Table~\ref{tab:kc_lpc} answers the question with measurements. Kiem's remark
is right about the method he meant, and for a reason he did not state: coding
raw ADC samples losslessly means keeping about four bits of noise per value that
his FX16 step discards (Section~\ref{sec:kr_lpcraw}). His remark does not carry
over to the LPC of this thesis, which predicts in the same domain and along the
same axis as DRHE. That LPC is level with DRHE everywhere it was measured: a
little ahead on his scene and on ColoRadar before the correction, a little
behind after it with his dictionary, and ahead again once the dictionary is
fitted to the data. The ranking on compression ratio is within a few per cent
and depends on the coder; what separates the two algorithms is structure ---
DRHE streams, LPC needs the whole frame --- and that is the basis of the choice
made in Section~\ref{sec:headtohead}.

\section{Detection metrics and the noise floor}
\label{sec:kc_det}

Kiem's FX16 row loses \SI{0.59}{\percent} of detections and adds
\SI{1.18}{\percent} spurious ones; this thesis's loses \SI{3.77}{\percent} and
adds \SI{1.11}{\percent}. Both measure the same thing --- the cost of quantising
to 16-bit fixed point, since every lossless coder shares FX16's metrics --- but on
different sensors, scenes and chains, and they are not comparable. Two
differences probably dominate. The ColoRadar codes use 9.8 of 16 bits, so the
quantisation step is coarse relative to the weakest returns; the quantisation
study shows that filling the container removes almost all of the error
(Table~\ref{tab:jr_lfr}). And the detection chain here integrates sixteen
receive channels over a Doppler FFT taken across transmitter-interleaved ramps
(Section~\ref{sec:refchain}), which spreads each target over several Doppler
cells and makes its local maxima more fragile.

The noise floors ($-70.7$ against $-105.4$\,dBFS) are not comparable at all:
0\,dBFS is not defined in Kiem's thesis, the sensors and gains differ, and NCI
over 16 channels and a 192-point Doppler FFT change the level. Within one
dataset, the noise floor is useful only for comparing schemes with each other.

\section{What the comparison finally says}
\label{sec:kc_summary}

\begin{keybox}
\textbf{Same:} the DRHE algorithm, to the bit; the finding that run-length
coding adds almost nothing; that DRHE beats every lossy scheme in the design
space on the combination of ratio and detection; that the design meets
\SI{100}{\mega\hertz} on his speed grade and fits his device.

\textbf{Different, and explained:} every hardware resource figure and the
initiation interval (float datapath, packer, closed loop, state clearing ---
each measured by the ladder); the throughput (a factor of two, once his
theoretical binary figure is compared with ours); the value of his fixed
dictionary on other data; the LPC ranking, which holds for his LPC and not for
ours.

\textbf{Not comparable:} the compression ratios and detection metrics on the
two datasets. The algorithms are the same; the datasets and their scaling are
not, and the similarity of 3.25 and 3.30 is a coincidence.

\textbf{Unexplained:} his published LUT and flip-flop counts exceed what his
described design needs by about 2.7 and 1.8 times (Section~\ref{sec:acc_lutff});
and his reported noise floor is 17\,dB higher than a white-noise scene with his
compression ratio would have (Section~\ref{sec:kr_nfcr}). The second most
likely reflects the difference between a real noise floor and a simulated one.
\end{keybox}
